\documentclass[10pt,a4paper]{amsart}
\usepackage[margin=19mm]{geometry}
\usepackage{amsmath,amssymb,mathtools}
\usepackage[scr=rsfs,cal=euler]{mathalfa}
\usepackage{xfrac}
\usepackage[unicode,hidelinks,bookmarks=false]{hyperref}
\newcommand{\ie}{i.e.\ }
\newcommand{\cf}{cf.\ }
\newcommand{\resp}{respectively\ }
\newcommand{\Subsection}{\subsection}
\providecommand{\nobreakdash}{\nobreak\hbox{-}}
\setlength{\emergencystretch}{2em}
\multlinegap0pt
\usepackage{amssymb}
\usepackage{xcolor}
\usepackage{algorithm}\usepackage{algpseudocode}
\newtheorem{lemma}{Lemma}
\title{On graphical partitions with restricted parts}
\author{Gilad Levy}
\date{Source: arXiv:2510.00007v3 (CC BY 4.0)}
\begin{document}
\maketitle

\noindent\emph{Source and licence.} On graphical partitions with restricted parts. Version arXiv:2510.00007v3 (CC BY 4.0), distributed by arXiv under the Creative Commons Attribution 4.0 International licence, http://creativecommons.org/licenses/by/4.0/, as recorded in the arXiv licence metadata for this exact version and shown on the abstract page https://arxiv.org/abs/2510.00007v3. Full licence notice: \texttt{licenses/arxiv-2510.00007-CC-BY-4.0.txt}.\par\smallskip

\noindent\textit{Editorial scope.} Counted unit: the article's lemma D-lemma (source0.tex lines 442--448). The article's complete original proof of it is delivered with it (source0.tex lines 451--483, one \texttt{proof} environment of 33 lines). No mathematics is written, repaired or re-derived: the statement and the proof are the article's own lines, in the article's own notation, and no retained line is substituted or re-flowed.  Retained uncounted as the source-local context the proof needs: lines 434--436.  Nothing was omitted from any retained span, so the package carries no omission record.  No retained line contains a citation, so the wrapper carries no bibliography and no bibliography entry.  No retained line contains a figure environment, an image-inclusion command or drawing code, so no image is delivered and the package declares no asset dependency.  Editorial formatting only, in addition: an AMS \texttt{amsart} preamble that restates the article's own declarations and macros used by the retained lines, namely the environments \texttt{lemma} on separate counters, together with the standard packages those lines need.  The retained environments print in this document as Lemma 1 in order of appearance, whereas the article prints the same items with the numbering of its own full document.  The complete native source of the article as distributed in the arXiv e-print for this exact version is bundled unchanged as \texttt{sources/186104/source0.tex}.
    \begin{equation}\label{beta}
        \frac{n-D^2}{2}=\sum_{\substack{m\in\mathbb{N} \\ \mu(m)\le D}}\frac{\mu(m)}{e^{\beta\mu(m)}-1}.
    \end{equation}

\begin{lemma}\label{D-lemma}
    Denote by $N(n,D,\mu)$ the number of partitions of $n$ from $\mathcal{P_{\mu}}$ with Durfee square side length $D$. Then
    \begin{equation}
        N(n,D,\mu)=e^{\beta(n-D^2)}\prod_{\substack{m\in\mathbb{N} \\ \mu(m)\le D}}\big(1-e^{-\beta\mu(m)}\big)^{-2} \cdot O(\beta^{3/2}),
    \end{equation}
    where $\beta$ is defined by~\eqref{beta}.
\end{lemma}

\begin{proof}
    Firstly, we denote by $N(n,D,\mu)$ the number of partitions of $n$ from $\mathcal{P_{\mu}}$ with a Durfee square side length of $D$. Andrews [12] showed that for unrestricted partitions, with $\mu(x)\equiv x$, it holds that
    \begin{equation*}
        N(n,D)=\big[q^{n-D^2}\big]\prod_{\substack{m=1}}^\infty(1-q^{m})^{-2},
    \end{equation*}
    where $q\in(0,1)$. Then by a natural generalization, for any $\mu\in M$ the following holds
    \begin{equation}
        N(n,D,\mu)=\big[q^{n-D^2}\big]\prod_{\substack{m\in\mathbb{N} \\ \mu(m)\le D}}\big(1-q^{\mu(m)}\big)^{-2}.
    \end{equation}
    Thus, we obtain
    \begin{equation*}
        N(n,D,\mu)=\frac{1}{2\pi i}\oint q^{-n+D^2-1}\prod_{\substack{m\in\mathbb{N} \\ \mu(m)\le D}}\big(1-q^{\mu(m)}\big)^{-2}dq
    \end{equation*}
    where the integration is taken along an infinitesimally small circular contour centered at the origin. Setting $q=e^{-\beta}$ we have
    \begin{equation*}
        N(n,D,\mu)=\frac{1}{2\pi i}\oint \exp\bigg[{-\beta(n-D^2)-\sum_{\substack{m\in\mathbb{N} \\ \mu(m)\le D}} 2\log\big(1-e^{-\beta \mu(m)}\big)}\bigg]d\beta.
    \end{equation*}
    This intgeral can be approximated using the saddle point method, in a similar manner as conducted above. We shall denote the logarithm of the integrand by $\Phi(\beta)$, and find its maximum point by differentiating:
    \begin{equation*}
        \Phi'(\beta)=n-D^2-\sum_{\substack{m\in\mathbb{N} \\ \mu(m)\le D}}\frac{2\mu(m)}{e^{\beta\mu(m)}-1}=0.
    \end{equation*}
    Hence, we obtained an equation for the parameter $\beta$, in terms of $n, D$ and $\mu$. In a similar behavior as for the parameter $\alpha$, we have $\beta\to0$ as $(n-D^2)\to\infty$, for any $\mu\in M$. \\
    \\
    Furthermore, the second derivative satisfies
    \begin{equation}\label{second-der}
        \Phi''(\beta)=\sum_{\substack{m\in\mathbb{N} \\ \mu(m)\le D}}\frac{2\mu^2(m)e^{\beta\mu(m)}}{(e^{\beta\mu(m)}-1)^2}=\int_0^D\frac{2x^2e^{\beta x}}{(e^{\beta x} -1)^2}\frac{dx}{\mu'(\mu^{-1}(x))} + O(\beta)=O\Big(\frac{1}{\beta^3}\Big).
    \end{equation}
    Thus, according to the saddle point method,
\begin{equation*}
    N(n,D,\mu)=\frac{e^{\Phi(\beta)}}{\sqrt{\Phi''(\beta)}}\cdot O(1)=e^{\beta(n-D^2)}\prod_{\substack{m\in\mathbb{N} \\ \mu(m)\le D}}\big(1-e^{-\beta\mu(m)}\big)^{-2} \cdot O(\beta^{3/2}),
\end{equation*}
and the proof is complete.
\end{proof}

\end{document}
